problem:  
Let \((M^n,g)\) be a smooth Riemannian manifold with \(\operatorname{Ric}_g \ge (n-1)k\,g\) for some real \(k\).  
As \(v\to0\), isoperimetric regions of volume \(v\) concentrate near points of maximal scalar curvature, and one has the second‑order expansion  
\[
I_M(v)=c_n v^{\frac{n-1}{n}}\Big(1 - A(p(v))v^{2/n}+o(v^{2/n})\Big),
\]
where \(A(p)\) is an explicit multiple of \(\mathrm{Scal}_g(p)\).

**Conjecture (Emergence of curvature invariants in higher‑order isoperimetric asymptotics):**  
For arbitrary \(n\ge3\) and any manifold with \(\operatorname{Ric}\ge (n-1)k\,g\), there exists a *third asymptotic term* in the small‑volume expansion  
\[
I_M(v)=c_n v^{\frac{n-1}{n}}\Big(1 - A(p)v^{2/n} + B(p)v^{4/n} + o(v^{4/n})\Big),
\]
where \(B(p)\) is expressible as a universal linear combination of scalar curvature derivatives and quadratic curvature invariants, e.g.  
\[
B(p)=\alpha_n |\mathrm{Ric}_g(p)|^2 + \beta_n |\mathrm{Weyl}_g(p)|^2 + \gamma_n \Delta\mathrm{Scal}_g(p).
\]
Characterize the tuple \((\alpha_n,\beta_n,\gamma_n)\) and determine whether its structure is dimension‑universal.  
In particular, does \(B(p)\) capture *all independent curvature invariants* appearing at cubic order in the metric’s normal‑coordinate expansion, or are there identities constraining these coefficients under the Ricci lower bound?

---

Why is it a "good" problem:  

1. **Mathematically sound and purely exploratory:**  
   Unlike rigidity statements relying on unestablished formulas, this conjecture asks first whether such higher‑order asymptotics exist and what precise curvature information they encode. This is a mathematically clean problem at the frontier of geometric analysis.

2. **Direct connection to the core of Riemannian geometry:**  
   Understanding \(B(p)\) would extend fundamental results linking curvature and isoperimetry, in the same way that the second‑order term revealed the role of the scalar curvature. A resolution would deepen our comprehension of how local curvature data influence global geometric inequalities.

3. **Cross‑disciplinary input:**  
   The problem sits at the intersection of geometric measure theory (variation of minimal hypersurfaces), Riemannian asymptotic expansions, and microlocal analysis. Progress would likely require new techniques uniting PDE expansions, curvature analysis, and minimal‑surface geometry.

4. **Potential to generate structural understanding:**  
   Knowing exactly which curvature contractions appear in \(B(p)\) would reveal how isoperimetric asymptotics filter geometric information: which invariants are “visible” to the isoperimetric profile, which are not, and how curvature bounds constrain the combinations allowed.

5. **Simple yet profound:**  
   At heart, the question asks: *What is the next curvature term that the isoperimetric profile “feels”?* This is easy to grasp yet demands deep computation and geometric insight, and any answer would add fundamentally new understanding to curvature‑dependent inequalities.  

Hence, this problem fulfills the qualities of a *good research problem*: it is clearly stated, foundational, logically sound, and bridges analysis and geometry toward uncovering previously unexplored structure in the relationship between curvature and isoperimetry.